\chapter[Introdução]{Introdução}
% \addcontentsline{toc}{chapter}{Introdução}

\section{Motivação}

O mercado de ações é reconhecido como um componente vital de qualquer economia, refletindo as tendências e o desempenho macroeconômico de um país. Índices como o IBOVESPA no Brasil e o S\&P 500 nos Estados Unidos servem como termômetros do mercado, orientando investidores e analistas \cite{Ferro2024AvaliacaoEnsemble}. A capacidade de antecipar seus movimentos pode otimizar estratégias de investimento, maximizar lucros e mitigar riscos. Com o avanço da Inteligência Artificial (IA), técnicas de aprendizado de máquina, em especial as redes neurais, têm se mostrado promissoras para modelar os padrões das séries temporais financeiras. No entanto, a previsão do comportamento desses índices ainda é um desafio, dada a natureza complexa, caótica, dinâmica e não-linear das séries temporais financeiras \cite{Ferro2024AvaliacaoEnsemble, Nelson2017UsoRedes}.

Apesar do debate sobre a previsibilidade dos mercados, o grande investimento em Aprendizado de Máquina, em particular, o \textit{Deep Learning}, tem aberto novos e promissores estudos para enfrentar este desafio. Os modelos de previsão baseados em ML e Redes Neurais (RN), como as Redes Neurais Recorrentes (RNN) e (LSTM), têm demonstrado a capacidade de capturar padrões e dependências temporais de longo prazo, superando métodos estatísticos tradicionais \cite{Nelson2017UsoRedes}.

Contudo, a sofisticação desses modelos, especialmente as redes neurais profundas, frequentemente resulta em um problema de ``caixa preta'' (\textit{black box}), onde o processo de tomada de decisão interna se torna opaco e de difícil interpretação. \cite{Ayatollahi2025XAIStockCrashes}. Essa falta de transparência compromete a confiança e a adoção prática dessas tecnologias. A falta de interpretabilidade impede que investidores e analistas compreendam quais fatores ou padrões de dados estão influenciando uma determinada previsão, tornando arriscado basear decisões financeiras em um sistema opaco. \cite{Gupta2025Interpretable}

Neste contexto, a Inteligência Artificial Explicável (\textit{xAI - Explainable AI}) surge como um campo crucial, buscando desenvolver métodos para tornar as decisões de modelos de IA compreensíveis para os seres humanos. Técnicas como SHAP são essenciais para atribuir a contribuição de cada variável à previsão do modelo, transformando a ``caixa preta'' em um sistema mais compreensível \cite{Ayatollahi2025XAIStockCrashes}. A aplicação de xAI no domínio financeiro é uma fronteira de pesquisa recente e de alta relevância, pois promete unir o poder preditivo das redes neurais com a transparência necessária para sua aplicação no mundo real.


\section{Contextualização do Problema}

A previsibilidade do mercado de ações é explorada por meio de diferentes técnicas. Modelos de \textit{Machine Learning} clássico, como Árvores de Decisão, e abordagens de \textit{Ensemble}, como \textit{Gradient Boosting}, são amplamente utilizados. Mais recentemente, as Redes Neurais Recorrentes (RNN), em especial a arquitetura LSTM, ganharam foco nos estudos de séries temporais, como abordaram os trabalhos \citeonline{Dametto2018EstudoRN}, \citeonline{Nelson2017UsoRedes} e \citeonline{Gupta2025Interpretable}.

A seleção da ferramenta mais adequada para previsão de séries temporais financeiras envolve tanto linguagens de programação de alto desempenho, como C++ e Java, quanto ambientes voltados à prototipação rápida, como Python. Em aplicações industriais é comum empregar bibliotecas otimizadas em C++ para a implementação dos modelos, enquanto na pesquisa acadêmica e no desenvolvimento experimental prevalecem plataformas de mais alto nível \cite{site9}. A escolha do \textit{framework} depende da natureza do problema, do modelo de Aprendizado de Máquina escolhido e do objetivo da previsão, que pode ser a estimativa de valores futuros (regressão), a classificação da direção do movimento do mercado ou a detecção de eventos raros (\textit{crash}) \cite{site6}.

A escolha dos dados de entrada é igualmente crucial, variando desde índices de bolsa (IBOVESPA, S\&P 500) até cotações de ações específicas, enriquecidos com indicadores de análise técnica. As séries temporais financeiras são caracterizadas por sua não-linearidade, volatilidade e estocasticidade, o que as torna intrinsecamente difíceis de prever. A modelagem eficaz exige que os algoritmos considerem, além dos dados históricos de preço, a influência de fatores externos (variáveis exógenas) como indicadores de análise técnica (AT) ou cotações de moedas (ex: o Dólar) \cite{Dametto2018EstudoRN}.

A evolução dos modelos de previsão mostra uma distinção fundamental entre arquiteturas simples e complexas. Enquanto os modelos mais básicos (análogos a redes de camada única ou \textit{single layer}) frequentemente apresentam limitações na captura de padrões complexos e dependências de longo prazo, arquiteturas como o \textit{Multilayer Perceptron} (MLP) e as redes de \textit{Deep Learning} (como LSTM e CNN) são capazes de modelar problemas não-linearmente separáveis graças à sua composição de múltiplas camadas intermediárias \cite{Ferro2024AvaliacaoEnsemble}.

Para superar as limitações dos modelos individuais, as técnicas de \textit{ensemble} foram desenvolvidas com o objetivo de combinar múltiplos modelos (homogêneos ou heterogêneos) para aumentar a robustez e acurácia da previsão. Essas abordagens incluem estratégias como \textit{Bagging}, \textit{Stacking} e Sistemas Híbridos Residuais \cite{Ferro2024AvaliacaoEnsemble}. O desempenho desses modelos é avaliado utilizando métricas de desempenho tradicionais de regressão, como o Erro Quadrático Médio (MSE), Raiz do Erro Quadrático Médio (RMSE), Erro Absoluto Médio (MAE) e Erro Percentual Absoluto Médio (MAPE). Em certas análises, uma métrica de custo-benefício é aplicada para verificar se o aumento na complexidade do modelo é justificado pelo ganho de desempenho e eficiência computacional \cite{Dametto2018EstudoRN}.

Apesar do aumento da precisão, modelos complexos como os ensembles e redes neurais profundas permanecem como incógnitas. O problema central abordado neste trabalho é a falta de interpretabilidade dos modelos de redes neurais aplicados à previsão do mercado de ações. Para enfrentar essa lacuna, este trabalho foca na aplicação de técnicas de xAI para elucidar o comportamento de um modelo de previsão de mercado. A questão de pesquisa que guia este estudo é: \textit{Como modelos baseados em LSTM, xLSTM e Transformer, integrados a técnicas de Inteligência Artificial Explicável (SHAP), podem melhorar a previsibilidade direcional e a interpretabilidade das decisões em séries temporais do índice S\&P500?}


\section{Impacto}

A capacidade de previsão de séries temporais financeiras é vital, pois fornece aos investidores informações de referência essenciais para avaliações de curto, médio e longo prazo, conferindo maior segurança ao investimento. Este trabalho contribui diretamente ao campo do \textit{algorithmic trading} ao focar em modelos de previsão de alto desempenho.

Em particular, a utilização de Redes Neurais como a LSTM é altamente relevante. As LSTMs são variações de redes recorrentes que se destacam por sua capacidade de reter dependências temporais de longo prazo, o que as torna modelos de \textit{Deep Learning} mais adequados para a modelagem da natureza sequencial e volátil dos dados do mercado de ações. A tecnologia de Redes Neurais que o trabalho aborda é, portanto, uma das mais promissoras no campo \cite{Gupta2025Interpretable}.

Embora o poder preditivo de modelos complexos como os \textit{ensembles} e LSTMs seja inegável, entender o porquê de uma previsão é tão ou mais importante do que a sua precisão. A abordagem deste trabalho, que integra modelos de previsão avançados com a Inteligência Artificial Explicável, ataca diretamente este problema. O uso de técnicas como SHAP é a relevância da tecnologia abordada, pois elas fornecem a interpretabilidade local e global necessária, transformando a opacidade em transparência. \cite{Benhamou2021XAIPlanning}

O impacto se manifesta por meio das seguintes contribuições:
\begin{enumerate}
    \item Desenvolver um modelo de previsão robusto para o mercado de ações que combine as capacidades de Redes Neurais com a interpretabilidade da xAI.
    \item Apresentar um \textit{framework} prático para a integração de xAI no processo de \textit{Deep Learning} aplicado à previsão de séries temporais financeiras.
    \item Oferecer \textit{insights} sobre quais fatores realmente influenciam as decisões do modelo (via SHAP).
\end{enumerate}

Dessa forma, o estudo contribui para o avanço de sistemas de análise financeira que não são apenas preditivamente robustos, mas também transparentes e confiáveis.

\section{Objetivos}

Este trabalho tem como propósito desenvolver, treinar e avaliar modelos de \textit{Deep Learning} baseados em LSTM, xLSTM e Transformer para previsão direcional do índice S\&P500, integrando técnicas de Inteligência Artificial Explicável (SHAP) a fim de tornar transparentes os fatores temporais e técnicos que influenciam as decisões do modelo. Para alcançar esse resultado, foram estabelecidos os seguintes objetivos:

\begin{itemize}
    \item Construir um pipeline de aquisição, pré-processamento e engenharia de atributos para séries temporais diárias do S\&P500.
    \item Implementar um modelo baseline LSTM para classificação direcional (sobe/desce) e realizar um ablation study sobre representações de entrada e estratégias de normalização.
    \item Implementar modelos xLSTM-TS e Transformer-TS para previsão do valor absoluto, de acordo com modelos de previsão validados, e compará-los ao baseline LSTM em termos de desempenho preditivo e custo computacional.
    \item Definir e aplicar métricas de avaliação para problemas desbalanceados de previsão direcional (MCC, Acurácia, F1, AUC-ROC) e, quando apropriado, métricas de regressão (RMSE, MAPE, R2).
    \item Integrar os modelos treinados ao framework SHAP para análise local e global das contribuições de atributos e \textit{timesteps}, produzindo explicações sobre o comportamento aprendido.
\end{itemize}

\section{Organização do Trabalho}

Este trabalho está organizado da seguinte forma. No capítulo 2 é desenvolvida a Revisão Bibliográfica, onde são apresentados os principais conceitos teóricos relacionados à previsão de valores de séries temporais no mercado financeiro, bem como os trabalhos correlatos que fundamentam a pesquisa. Em seguida, no capítulo 3, na Metodologia, estão descritos os procedimentos adotados para a realização do estudo, incluindo as estratégias de coleta e tratamento dos dados, as ferramentas e técnicas utilizadas, além dos critérios de avaliação empregados. Os Capítulos 4 e 5 apresentam os resultados e as suas interpretações, tanto em termos de desempenho preditivo quanto de interpretabilidade. Por fim, o Capítulo 6 discute as conclusões do trabalho, suas limitações e sugestões para pesquisas futuras.
